\chapter{Lua API reference}
This document describes the Lua runtime model used by the firmware and the currently exposed Lua namespaces (modules) implemented in C.

The goal is to help you write two Lua files:

\begin{itemize}
\item App script (app.lua) with UI logic (drawing on display, handling loop, etc.).
\item Sensor script (sensor/*.lua) that provides measured values through sensor.get(index).
\end{itemize}
\section{Runtime model}
\subsection{App contract}
Each app file returns a table (commonly M) with callbacks:

\begin{itemize}
\item M.loop() required
\item M.init() optional
\item M.event(event\_type) optional
\item M.cleanup() optional
\end{itemize}
The firmware loads your app, stores these callbacks, and executes loop() repeatedly. If loop is missing, the app is rejected. On runtime errors, firmware can switch to a fallback app.

\subsection{Sensor contract}
Sensor Lua file should return a table named sensor with:

\begin{itemize}
\item sensor.get(index) returning a numeric value.
\end{itemize}
The firmware loads the sensor script and stores it globally as sensor, then app scripts call sensor.get(...).

\subsection{Globals provided by firmware}
\begin{itemize}
\item apppath — path to current app folder (useful for loading assets/fonts relative to app location).
\item sensor — table returned by active sensor script.
\end{itemize}
\section{Lua namespaces (modules)}
The firmware registers the following Lua modules:

\begin{itemize}
\item display
\item system
\item onewire
\item nvs
\item ble
\end{itemize}
Load them in Lua with require("module") (or via pcall(require, "module") for optional usage).

\section{API reference}
\section{display}
\subsection{Resource management}
\begin{itemize}
\item display.loadimg(path) -> img\_handle|nil
\end{itemize}
\begin{itemize}
\item Decodes image and returns handle (integer-like pointer).
\end{itemize}
\begin{itemize}
\item display.loadfont(path) -> font\_handle|nil
\item display.free(handle)
\end{itemize}
\begin{itemize}
\item Frees image/font allocated by loadimg / loadfont.
\end{itemize}
\subsection{Frame control}
\begin{itemize}
\item display.fill(color)
\end{itemize}
\begin{itemize}
\item Fills entire screen.
\end{itemize}
\begin{itemize}
\item display.update()
\end{itemize}
\begin{itemize}
\item Flushes frame and clears screen to black afterwards.
\end{itemize}
\subsection{Drawing primitives}
\begin{itemize}
\item display.drawimg(img, x, y)
\item display.string(font, x, y, text, color)
\item display.fillrect(x, y, width, height, color)
\item display.drawline(x1, y1, x2, y2, color)
\item display.drawcircle(x, y, r, color)
\item display.drawfillcircle(x, y, r, color)
\item display.drawroundrect(x1, y1, x2, y2, r, color)
\item display.drawarrow(x0, y0, x1, y1, w, color)
\item display.drawfillarrow(x0, y0, x1, y1, w, color)
\item display.drawpolygon(\{x1, y1, x2, y2, ...\}, color)
\end{itemize}
\begin{itemize}
\item Even count of coordinates required (at least 4 elements).
\end{itemize}
\subsection{Color constants}
\begin{itemize}
\item display.RED
\item display.GREEN
\item display.BLUE
\item display.WHITE
\item display.BLACK
\item display.YELLOW
\item display.CYAN
\item display.PURPLE
\item display.GRAY
\end{itemize}
\section{system}
\begin{itemize}
\item system.delay(ms)
\item system.backlight(level)
\item system.V2\_pullup(enable)
\item system.out5V(enable)
\item system.keys() -> int
\item system.adc\_ch(channel) -> number
\item system.get\_time() -> hour, minute, second
\item system.config\_sleep(enable)
\end{itemize}
\section{onewire}
\begin{itemize}
\item onewire.scan18b20() -> count
\item onewire.read18b20(index) -> temperature\_c
\end{itemize}
\section{nvs}
Namespace is backed by NVS partition namespace app.

\begin{itemize}
\item nvs.get(key) -> string|nil
\item nvs.push(key, value\_string) -> boolean
\end{itemize}
\section{ble}
\begin{itemize}
\item ble.init() -> true
\item ble.get\_telemetry() -> rpm, speed, gear
\end{itemize}
\section{Example \#1 — Simple voltmeter app}
The coordinates below fit the 160 by 80 pixel screen, and the bar width is clamped.
Below is a minimal app using the same style as existing voltmeter apps (sensor.get(1), font draw, cleanup).

\begin{verbatim}
local M = {}
local font = display.loadfont(apppath .. "/../fonts/H24XB.FNT")

local function read_voltage()
return sensor.get(1)
end

function M.loop()
local v = read_voltage()
display.fill(display.BLACK)
display.string(font, 10, 30,
    string.format("%2.1f V", v), display.WHITE)
display.fillrect(10, 60,
    math.max(0, math.min(140, math.floor((v - 9.5) * 12))),
    8, display.GREEN)
display.update()
end

function M.cleanup()
display.free(font)
end

return M
\end{verbatim}
Notes:

\begin{itemize}
\item display.update() is called once per frame.
\item cleanup() should free loaded resources.
\item Clamp bar calculations in production code (existing apps clamp values to display range).
\end{itemize}
\section{Example \#2 — Sensor Lua script (NTC-like)}
This script demonstrates reading ADC, enabling pull-up rail, smoothing, and returning processed value.

\begin{verbatim}
sensor = {}

local filt = 0.0

function sensor.get(i)
i = i or 0
system.V2_pullup(1)

local raw = system.adc_ch(i)
filt = filt + (raw - filt) / 10.0

-- Example transform (replace with your real calibration/math):
local value = filt * 1.0
return value
end

return sensor
\end{verbatim}
Notes:

\begin{itemize}
\item Keep sensor.get(i) fast; it is called from app loops.
\item You can store calibration in NVS (for example, via nvs.get("key") and tonumber(...)).
\end{itemize}
\section{Practical tips}
\begin{itemize}
\item Use pcall(require, "nvs") when NVS access is optional.
\item Build paths from apppath for portable assets (apppath .. "image.jpg", apppath .. "/../fonts/...").
\item Always free display.loadimg / display.loadfont handles in cleanup().
\item If using DS18B20, call onewire.scan18b20() once and onewire.read18b20(index) in loop.
\end{itemize}
